% Insert before Discussion/Conclusion. Uses only the original template packages.
\section{Conditional GAN for Ban/Pick Simulation}
\label{sec:bpgan}

\subsection{Task and Scope}
To extend the project from outcome prediction to generation, a conditional
generative adversarial network (cGAN) was trained to generate a lineup given
a map, a round category, and supplied character bans. A generated lineup
contains one hunter and four distinct survivors. This is a ban-conditioned
lineup-completion task, not a complete sequential draft policy: bans are
provided as inputs rather than generated as actions.

The objective is to approximate observed selection patterns, not to maximize
winning probability. No win probability is assigned to the generated lineups.
The current decoder enforces faction membership, supplied bans, and survivor
uniqueness, but does not implement all patch-specific or tournament-specific
rules. The method follows the conditional adversarial framework of Mirza and
Osindero\footnote{\url{https://arxiv.org/abs/1411.1784}} and uses a
Gumbel-Softmax relaxation for categorical
generation\footnote{\url{https://arxiv.org/abs/1611.01144}}.

\subsection{Data and Chronological Split}
Starting from 5,988 single-game records, preprocessing excludes 18 non-regular
games, five incomplete or duplicate-pick records, 14 records without a map,
and 14 records whose picks conflict with the supplied bans. Conditions use
the initial active-ban slots and recorded global bans; later active-ban slots,
scores, and outcome labels are excluded. These fields are treated as a
context snapshot, not a verified chronological action log.

Training uses 4,952 cleaned games from 2020 Summer through 2024 Autumn.
Validation uses 2025 Summer, and testing uses 2025 Autumn. Full matches do not
cross split boundaries. Character and map vocabularies are fitted on training
data only, resulting in 26 hunter characters and 45 survivor characters.

Only games whose selected characters and maps occur in the training vocabulary
are supported. This retains 177 of 501 cleaned validation games (35.3\%) and
106 of 484 cleaned test games (21.9\%). This low coverage is a central
limitation: the results describe the supported subset and must not be
generalized to the entire test season or to unseen characters.

\subsection{Architecture and Optimization}
The generator receives a 16-dimensional Gaussian noise vector $z$ and a
condition vector $c$ consisting of map and round indicators and binary ban
masks. Two fully connected hidden layers of 64 units with ReLU activations
produce five categorical heads: one hunter head and four survivor heads.
Straight-through Gumbel-Softmax samples allow adversarial gradients to reach
the generator. During decoding, banned characters are masked and each selected
survivor is removed from subsequent survivor heads.

The discriminator has two 64-unit hidden layers with leaky ReLU activations.
It receives the condition, the hunter indicator, and the sum of survivor
indicators divided by four. This set representation prevents arbitrary
survivor slot order from becoming a real-versus-generated shortcut.
Its objective is
\begin{equation}
L_D=-\mathbb{E}_{c,y}\log D(c,y)
    -\mathbb{E}_{c,z}\log[1-D(c,G(c,z))].
\end{equation}
The generator uses the non-saturating objective
\begin{equation}
L_G=-\mathbb{E}_{c,z}\log D(c,G(c,z))
    +\alpha L_{\mathrm{CE}},
\end{equation}
where $L_{\mathrm{CE}}$ is the mean supervised cross-entropy across five
categorical heads, with survivor target order randomized during training.
The auxiliary term encourages observed pick frequencies rather than a
particular survivor ordering.

Two candidates, $\alpha=0$ and $\alpha=1$, were each trained for 60 epochs
using Adam with learning rate 0.0003, batch size 128, and
$(\beta_1,\beta_2)=(0.5,0.9)$. The random seed was 20261003.
Checkpoints were evaluated at epoch 1 and every five epochs thereafter.
The hybrid candidate ($\alpha=1$) at epoch 40 was selected by validation
divergence, without refitting or selecting on test metrics.

\subsection{Evaluation and Preliminary Results}
For each supported context, eight lineups are sampled. Distribution similarity
is measured by the mean of hunter and survivor marginal Jensen--Shannon
divergences, using natural logarithms:
\begin{equation}
\mathrm{JS}(P,Q)=\frac{1}{2}\mathrm{KL}(P\Vert M)
               +\frac{1}{2}\mathrm{KL}(Q\Vert M),
\qquad M=(P+Q)/2.
\end{equation}
This measures marginal pick-frequency agreement, not conditional tactical
quality or joint lineup quality. A baseline samples add-one-smoothed training
map frequencies and applies the same decoding constraints.

\begin{table}[htbp]
\centering
\caption{Generation on 106 supported test contexts, with 848 samples per method.
Lower JS is better. Diversity and novelty are not win-rate measures.}
\label{tab:gan}
\begin{tabular}{lrrr}
\toprule
Method & Mean JS & Diversity (\%) & Novelty (\%) \\
\midrule
Map-frequency sampling & 0.2007 & 100.0 & 96.7 \\
Hybrid cGAN & 0.0972 & 99.8 & 94.0 \\
\bottomrule
\end{tabular}
\end{table}

Diversity is the number of unique unordered lineups divided by eight, averaged
across contexts. Novelty is the fraction of lineups absent from training,
ignoring context. High diversity or novelty alone does not establish useful
drafting behavior. Without decoding constraints, 77.1\% of cGAN samples and
63.8\% of baseline samples violate at least one supplied-ban or uniqueness
constraint. With constraints, both violation rates are zero. This zero rate
is guaranteed by the decoder and is not evidence that the GAN independently
learned tournament legality.

The hybrid model improves marginal distribution similarity on the supported
test subset, but the experiment uses only one seed and excludes most test
games because of vocabulary limitations. Broader evaluation should include
multiple seeds, conditional and joint-distribution checks, expert assessment,
and explicit season rules. Connecting generated lineups to an outcome model
would require separate validation and would still not establish causal
advantages from changing a draft.

\subsection{Interactive Demonstration}
The trained generator weights are exported to JSON and evaluated directly in
the project website. Users can select a map and round, edit banned characters,
and supply a reproducible random seed. The interface generates alternative
lineups and checks the implemented constraints. It deliberately reports no
predicted win-rate improvement or claim of optimal play.
